% ID: SEL-D0064-OA-E01
% ESTADO: cotejada. Segunda pasada visual y borrador compilado revisados.
% Fuente: Selectividad/sel_2010_matematicas_aplicadas/mat_apli/mat_apli_sept/mat_apli_sept/EXAMEN-ANDALUCÍA-3.pdf
% SHA-256: 9c3d2d628115a2f20f17b47171106ddd90f214b5f6804bdf9a7dcb7037c75efd
% Página: 1; opción A; ejercicio original 1.

Sea el recinto del plano definido por el siguiente sistema de inecuaciones:
\[
3x+y\geq4;\quad x+y\leq6;\quad0\leq y\leq5.
\]
\begin{enumerate}[label=\alph*)]
\item \textbf{(1 punto)} Represéntelo gráficamente.
\item \textbf{(1 punto)} Calcule los vértices de dicho recinto.
\item \textbf{(0.5 puntos)} En el recinto anterior, halle los valores máximo y mínimo de la función $F(x,y)=5x+3y$. \textquestiondown{}En qué puntos se alcanzan dichos valores?
\end{enumerate}
